\documentclass[11pt,a4paper]{article}
\usepackage[top=2.3cm,bottom=2.3cm,left=2.2cm,right=2.2cm]{geometry}

\usepackage{amsmath,amssymb}
\usepackage{fontspec}
\setmainfont{Liberation Serif}
\setsansfont{Liberation Sans}
\setmonofont{DejaVu Sans Mono}[Scale=0.84]

\usepackage{xcolor}
\definecolor{brandblue}{RGB}{20, 55, 105}
\definecolor{accentblue}{RGB}{35, 95, 165}
\definecolor{codebg}{RGB}{248, 249, 251}
\definecolor{codeframe}{RGB}{215, 222, 230}
\definecolor{javakeyword}{RGB}{0, 45, 170}
\definecolor{javastring}{RGB}{5, 120, 20}
\definecolor{javacomment}{RGB}{120, 120, 120}
\definecolor{javanumber}{RGB}{20, 75, 220}

\definecolor{noteborder}{RGB}{30, 110, 65}
\definecolor{notebg}{RGB}{242, 249, 245}
\definecolor{warnborder}{RGB}{185, 75, 10}
\definecolor{warnbg}{RGB}{254, 248, 240}
\definecolor{tipborder}{RGB}{30, 95, 160}
\definecolor{tipbg}{RGB}{242, 247, 254}

\usepackage{fancyhdr}
\usepackage{lastpage}
\pagestyle{fancy}
\fancyhf{}
\lhead{\parbox[t]{0.58\textwidth}{\raggedright\small\sffamily\color{brandblue}\textbf{T04 --- Sistema di Tipi, Primitive e Reference}}}
\rhead{\small\sffamily\color{gray}Programmazione II -- UniVR (2025/2026)}
\cfoot{\small\sffamily Pagina \thepage\ di \pageref{LastPage}}
\renewcommand{\headrulewidth}{0.5pt}
\renewcommand{\headrule}{\hbox to\headwidth{\color{brandblue!70!black}\leaders\hrule height \headrulewidth\hfill}}

\usepackage{titlesec}
\titleformat{\section}{\Large\bfseries\sffamily\color{brandblue}}{\thesection}{0.8em}{}[\titlerule]
\titleformat{\subsection}{\large\bfseries\sffamily\color{accentblue}}{\thesubsection}{0.8em}{}
\titleformat{\subsubsection}{\normalsize\bfseries\sffamily\color{brandblue!85!black}}{\thesubsubsection}{0.8em}{}

\usepackage{needspace}
\usepackage{fancyvrb}
\DefineVerbatimEnvironment{asciibox}{Verbatim}{
  fontsize=\fontsize{7.5pt}{9.2pt}\selectfont,
  frame=single,
  rulecolor=\color{codeframe},
  fillcolor=\color{codebg},
  framesep=2.5mm
}
\usepackage{listings}
\lstset{
  backgroundcolor=\color{codebg},
  frame=single,
  rulecolor=\color{codeframe},
  basicstyle=\ttfamily\footnotesize,
  keywordstyle=\color{javakeyword}\bfseries,
  stringstyle=\color{javastring},
  commentstyle=\color{javacomment}\itshape,
  numberstyle=\tiny\color{gray},
  numbers=left,
  numbersep=7pt,
  stepnumber=1,
  breaklines=true,
  breakatwhitespace=true,
  showstringspaces=false,
  tabsize=4,
  xleftmargin=12pt,
  framexleftmargin=12pt
}

\newcommand{\passthrough}[1]{#1}
\providecommand{\tightlist}{\setlength{\itemsep}{0pt}\setlength{\parskip}{0pt}}

\usepackage{tcolorbox}
\tcbuselibrary{skins,breakable}
\newtcolorbox{studynote}[1][Nota]{
  enhanced,
  breakable,
  colback=notebg,
  colframe=noteborder,
  fonttitle=\bfseries\sffamily\small,
  coltitle=white,
  title={#1},
  arc=2mm,
  boxrule=0.8pt,
  left=9pt, right=9pt, top=7pt, bottom=7pt
}

\newtcolorbox{studywarning}[1][Attenzione]{
  enhanced,
  breakable,
  colback=warnbg,
  colframe=warnborder,
  fonttitle=\bfseries\sffamily\small,
  coltitle=white,
  title={#1},
  arc=2mm,
  boxrule=0.8pt,
  left=9pt, right=9pt, top=7pt, bottom=7pt
}

\newtcolorbox{studytip}[1][Suggerimento]{
  enhanced,
  breakable,
  colback=tipbg,
  colframe=tipborder,
  fonttitle=\bfseries\sffamily\small,
  coltitle=white,
  title={#1},
  arc=2mm,
  boxrule=0.8pt,
  left=9pt, right=9pt, top=7pt, bottom=7pt
}

\usepackage{booktabs}
\usepackage{longtable}
\usepackage{array}
\usepackage{calc}

\usepackage{hyperref}
\hypersetup{
  colorlinks=true,
  linkcolor=accentblue,
  urlcolor=accentblue,
  citecolor=brandblue,
  pdfborder={0 0 0}
}

\title{\Huge\bfseries\sffamily\color{brandblue} T04 --- Sistema di Tipi, Primitive e Reference \\[0.3em] \large\sffamily Tipi Primitivi vs Reference, Casting, Wrapper, Operatori e Analisi di Example.java}
\author{\textbf{Docente:} Prof. Michele Pasqua \quad---\quad \textbf{Appunti Revisione:} Wally}
\date{A.A. 2025/2026 -- Universit\`a degli Studi di Verona}

\begin{document}
\maketitle
\thispagestyle{fancy}
\vspace{-1.2em}
\noindent\textcolor{brandblue}{\rule{\textwidth}{0.8pt}}
\vspace{1.2em}

Con questo blocco entriamo nel vivo della semantica del linguaggio:
analizziamo come la JVM gestisce i tipi, la memoria (Stack vs Heap), i
puntatori/reference, gli operatori logici e lo scope, confrontando la
teoria di
\href{file:///home/wally/Documenti/GitHub/Programmazione-II/25-26/Teoria-e-Esercitazioni/Slides-20260902/T04\%20-\%20Types\%20and\%20Objects.pdf}{T04
- Types and Objects.pdf} con il primo file di codice ufficiale del
docente:
\href{file:///home/wally/Documenti/GitHub/Programmazione-II/25-26/Teoria-e-Esercitazioni/Codice-20260902/Lezione04/Example.java}{\passthrough{\lstinline!Example.java!}}.

\begin{center}\rule{0.5\linewidth}{0.5pt}\end{center}

\subsubsection{1. Teoria: Concetti Teorici dalle
Slide}\label{teoria-concetti-teorici-dalle-slide}

\paragraph{A. Tipizzazione dei Linguaggi e il ``Contratto'' di Tipo
(Slide
1--7)}\label{a.-tipizzazione-dei-linguaggi-e-il-contratto-di-tipo-slide-17}

\begin{itemize}
\tightlist
\item
  \textbf{Definizione formale informale di Tipo}: Un tipo è una coppia
  formata da un \textbf{insieme di valori ammissibili} e da un
  \textbf{insieme di operazioni consentite} su di essi.

  \begin{itemize}
  \tightlist
  \item
    Es. \passthrough{\lstinline!int!}: valori da \(-2^{31}\) a
    \(2^{31}-1\), operazioni \passthrough{\lstinline!+!},
    \passthrough{\lstinline!-!}, \passthrough{\lstinline!*!},
    \passthrough{\lstinline!/!}, \passthrough{\lstinline!\%!}.
  \item
    Il tipo definisce un \textbf{contratto}: vincola cosa una variabile
    può contenere e in quali espressioni può comparire.
  \end{itemize}
\item
  \textbf{Weakly Typed (es. Python, JavaScript) vs Strongly Typed
  (Java)}:

  \begin{itemize}
  \tightlist
  \item
    Nei linguaggi \emph{weakly typed}, le variabili sono contenitori
    generici slegati dal tipo: una variabile \passthrough{\lstinline!x!}
    può contenere un intero, poi una stringa, poi una tupla.
  \item
    Nei linguaggi \emph{strongly typed} come Java, ogni variabile
    richiede una dichiarazione esplicita e immutabile di tipo a tempo di
    compilazione (\passthrough{\lstinline!int x;!}). Assegnare un tipo
    incompatibile genera un \textbf{Compile-Time Error}.
  \item
    \emph{Meme di JavaScript a lezione}:

    \begin{itemize}
    \tightlist
    \item
      \passthrough{\lstinline!"11" + 1!} \(\implies\)
      \passthrough{\lstinline!"111"!} (coercizione implicita a stringa e
      concatenazione).
    \item
      \passthrough{\lstinline!"11" - 1!} \(\implies\)
      \passthrough{\lstinline!10!} (coercizione implicita a numero e
      sottrazione numerica).
    \item
      In Java queste ambiguità ed errori silenziosi sono impossibili: il
      compilatore blocca ogni operazione non conforme prima
      dell'esecuzione.
    \end{itemize}
  \end{itemize}
\end{itemize}

\begin{center}\rule{0.5\linewidth}{0.5pt}\end{center}

\paragraph{B. I Tipi Primitivi in Java e le Conversioni (Slide
8--13)}\label{b.-i-tipi-primitivi-in-java-e-le-conversioni-slide-813}

Java dispone di \textbf{8 tipi primitivi} (allocati direttamente per
valore nello Stack o inline negli oggetti, non sono oggetti):

\begin{longtable}[]{@{}
  >{\raggedright\arraybackslash}p{(\linewidth - 6\tabcolsep) * \real{0.2500}}
  >{\raggedright\arraybackslash}p{(\linewidth - 6\tabcolsep) * \real{0.2500}}
  >{\raggedright\arraybackslash}p{(\linewidth - 6\tabcolsep) * \real{0.2500}}
  >{\raggedright\arraybackslash}p{(\linewidth - 6\tabcolsep) * \real{0.2500}}@{}}
\toprule\noalign{}
\begin{minipage}[b]{\linewidth}\raggedright
Tipo Primitivo
\end{minipage} & \begin{minipage}[b]{\linewidth}\raggedright
Dimensione
\end{minipage} & \begin{minipage}[b]{\linewidth}\raggedright
Range di Valori
\end{minipage} & \begin{minipage}[b]{\linewidth}\raggedright
Default (campi)
\end{minipage} \\
\midrule\noalign{}
\endhead
\bottomrule\noalign{}
\endlastfoot
\passthrough{\lstinline!byte!} & 8 bit (1 byte) & da \(-128\) a \(+127\)
(\(-2^7 \dots 2^7-1\)) & \passthrough{\lstinline!0!} \\
\passthrough{\lstinline!short!} & 16 bit (2 byte) & da \(-32.768\) a
\(+32.767\) (\(-2^{15} \dots 2^{15}-1\)) &
\passthrough{\lstinline!0!} \\
\textbf{\passthrough{\lstinline!int!}} & 32 bit (4 byte) & da
\(-2.147.483.648\) a \(+2.147.483.647\) (\(-2^{31} \dots 2^{31}-1\)) &
\passthrough{\lstinline!0!} \\
\passthrough{\lstinline!long!} & 64 bit (8 byte) & da \(-2^{63}\) a
\(+2^{63}-1\) (suffisso \passthrough{\lstinline!L!} o
\passthrough{\lstinline!l!}) & \passthrough{\lstinline!0L!} \\
\passthrough{\lstinline!float!} & 32 bit IEEE 754 & Precisione singola
(suffisso obbligatorio \passthrough{\lstinline!f!} o
\passthrough{\lstinline!F!}) & \passthrough{\lstinline!0.0f!} \\
\textbf{\passthrough{\lstinline!double!}} & 64 bit IEEE 754 & Precisione
doppia (default per numeri con la virgola) &
\passthrough{\lstinline!0.0d!} \\
\textbf{\passthrough{\lstinline!char!}} & 16 bit (2 byte) & da
\passthrough{\lstinline!0!} a \passthrough{\lstinline!65.535!} (Unicode
UTF-16, \textbf{unsigned}) & \passthrough{\lstinline!'\\u0000'!}
(NUL) \\
\textbf{\passthrough{\lstinline!boolean!}} & 1 bit logico &
\passthrough{\lstinline!true!} oppure \passthrough{\lstinline!false!} &
\passthrough{\lstinline!false!} \\
\end{longtable}

\begin{itemize}
\tightlist
\item
  \textbf{Aritmetica Modulare e Overflow degli
  \passthrough{\lstinline!int!}}:

  \begin{itemize}
  \tightlist
  \item
    In Java non viene lanciata alcuna eccezione se un calcolo intero
    supera il limite massimo: avviene il \textbf{wrap-around} secondo
    l'aritmetica del complemento a due a 32 bit.
  \item
    Se \passthrough{\lstinline!a = 2147483647!} (\(2^{31}-1\)),
    l'operazione \passthrough{\lstinline!a + 1!} restituisce
    \(-2147483648\) (\(-2^{31}\)).
  \end{itemize}
\item
  \textbf{Conversioni di Tipo (Casting)}:

  \begin{itemize}
  \tightlist
  \item
    \textbf{Widening (Allargamento / Implicito)}: da un tipo più stretto
    a uno più ampio (nessuna perdita d'ordine di grandezza):
    \[\text{byte} \longrightarrow \text{short} \longrightarrow \text{int} \longrightarrow \text{long} \longrightarrow \text{float} \longrightarrow \text{double}\]
    Es: \passthrough{\lstinline!float f = 3;!} converte automaticamente
    \passthrough{\lstinline!int 3!} in \passthrough{\lstinline!3.0f!}.
  \item
    \textbf{Narrowing (Restringimento / Esplicito obbligatorio)}: da un
    tipo più ampio a uno più stretto, richiede il cast sintattico
    \passthrough{\lstinline!(tipo)!} perché comporta potenziale
    troncamento o perdita di informazione: Es:
    \passthrough{\lstinline!int i = (int) 3.2;!} scarta la parte
    decimale e memorizza \passthrough{\lstinline!3!}.
  \end{itemize}
\item
  \textbf{I Caratteri sono Numeri Interi Unsigned a 16 bit}:

  \begin{itemize}
  \tightlist
  \item
    Un \passthrough{\lstinline!char!} memorizza il codice numerico
    Unicode (UTF-16 code unit).
  \item
    \passthrough{\lstinline!'V'!}, \passthrough{\lstinline!(char) 86!} e
    \passthrough{\lstinline!'\\u0056'!} sono tre rappresentazioni
    letterali dello \textbf{stesso identico dato binario} (valore
    decimale 86).
  \end{itemize}
\item
  \textbf{Tipi Non-Primitivi (Reference Types) e la Classe
  \passthrough{\lstinline!String!}}:

  \begin{itemize}
  \tightlist
  \item
    Tutto ciò che non è un tipo primitivo è un \textbf{oggetto}.
  \item
    \passthrough{\lstinline!String!} è una classe immutabile.
    L'operatore \passthrough{\lstinline!+!} tra una stringa e qualsiasi
    altro tipo (primitivo o oggetto) applica automaticamente
    l'overloading di concatenazione, invocando internamente la
    conversione a stringa.
  \end{itemize}
\end{itemize}

\begin{center}\rule{0.5\linewidth}{0.5pt}\end{center}

\paragraph{C. Modello di Memoria, Puntatori e Garbage Collection (Slide
14--21)}\label{c.-modello-di-memoria-puntatori-e-garbage-collection-slide-1421}

\begin{itemize}
\tightlist
\item
  \textbf{Assenza di Variabili Globali}: In Java non esiste lo scope
  globale in stile C. Ogni dato o funzione appartiene a una classe o a
  un'istanza.
\item
  \textbf{La Verità sui Puntatori in Java}:

  \begin{itemize}
  \tightlist
  \item
    \emph{«Java does have pointers, but there is no pointer
    arithmetic»}: ogni variabile di tipo non primitivo è tecnicamente un
    \textbf{puntatore/riferimento} (\emph{reference}) a un blocco di
    memoria nello \textbf{Heap}.
  \item
    La JVM protegge la memoria: il programmatore non può vedere
    l'indirizzo fisico, né fare \passthrough{\lstinline!ptr++!} o
    dereferenziare offset arbitrari.
  \item
    Il passaggio dei parametri ai metodi in Java è \textbf{sempre
    rigorosamente per valore} (\emph{pass-by-value}): per gli oggetti,
    viene copiato per valore il \textbf{riferimento} (l'indirizzo
    logico).
  \end{itemize}
\item
  \textbf{Rappresentazione di un Reference}:

  \begin{itemize}
  \tightlist
  \item
    Una variabile d'istanza non inizializzata contiene
    \passthrough{\lstinline!null!} (es.
    \passthrough{\lstinline!Vehicle@null!}).
  \item
    Quando viene invocato \passthrough{\lstinline!new Vehicle()!}, la
    JVM alloca l'oggetto nello Heap e restituisce l'indirizzo:
    \passthrough{\lstinline!myCar!} conterrà un valore del tipo
    \passthrough{\lstinline!Vehicle@25c7f37d!}.
  \end{itemize}
\item
  \textbf{Errori di Memoria e
  \passthrough{\lstinline!NullPointerException!} (NPE)}:

  \begin{itemize}
  \tightlist
  \item
    Se si tenta di invocare un metodo o accedere a un campo su una
    reference che vale \passthrough{\lstinline!null!} (es.
    \passthrough{\lstinline!myCar2.setFrameNumber(5)!}), la JVM solleva
    una \textbf{\passthrough{\lstinline!NullPointerException!}} a
    runtime, interrompendo l'esecuzione.
  \end{itemize}
\item
  \textbf{Perché i programmatori scrivono
  \passthrough{\lstinline!myCar = null;!}?}:

  \begin{itemize}
  \tightlist
  \item
    L'assegnamento a \passthrough{\lstinline!null!} rimuove il puntatore
    verso l'oggetto nello Heap. Se nessun'altra reference punta a
    quell'area di memoria, l'oggetto diventa \textbf{irraggiungibile}
    (\emph{unreachable}) ed eleggibile per la \textbf{Garbage Collection
    (GC)}, consentendo alla JVM di bonificare la memoria.
  \end{itemize}
\end{itemize}

\begin{center}\rule{0.5\linewidth}{0.5pt}\end{center}

\subsubsection{\texorpdfstring{2. Codice: Analisi Dettagliata del
Codice:
\href{file:///home/wally/Documenti/GitHub/Programmazione-II/25-26/Teoria-e-Esercitazioni/Codice-20260902/Lezione04/Example.java}{\texttt{Example.java}}}{2. Codice: Analisi Dettagliata del Codice: Example.java}}\label{codice-analisi-dettagliata-del-codice-example.java}

Analizziamo riga per riga il sorgente della Lezione 04, evidenziando
tutte le trappole e i dettagli d'esame:

\begin{lstlisting}[language=Java]
public class Example {
    static float f1;
    static float f2 = 3.4f;
    static String s1;
    static final double PI = 3.14; 
    
    public static void main(String[] args) {
        int i1 = 6;
        int i2;
        char c1 = 'V';
        char c2 = 86;
        char c3 = '\u0056';
        String s2;
            
        System.out.println("i1: " + i1);
        // System.out.println("i2: " + i2); // compile-time error
        System.out.println("f1: " + f1);
        System.out.println("f2: " + f2);
        System.out.println("c1: " + c1 + " c2: " + c2 + " c3: " + c3);
        System.out.println("s1: " + s1);
        /*
         System.out.println("s2: " + s2); // compile-time error
         System.out.println(s1.toString()); // run-time error
         PI = 3.0; // compile-time error
        */
        int k = 0;
        System.out.println(true | k++ == 0);
        System.out.println("k (standard eval): " + k);
        k = 0;
        System.out.println(true || k++ == 0);
        System.out.println("k (short-circuit eval): " + k);
        //
        int n = 4;
        n ++;
        // int n = 8; // compile-time error
        {
            int m = 7;
            m ++;
        }
        {
            int m = 8;
            m --;
        }
        System.out.println("n: " + n);
        // System.out.println("m: " + m); // compile-time error
    }
}
\end{lstlisting}

\paragraph{Analisi: Le Finezze Implementative Spiegate Riga per
Riga:}\label{analisi-le-finezze-implementative-spiegate-riga-per-riga}

\begin{enumerate}
\def\labelenumi{\arabic{enumi}.}
\tightlist
\item
  \textbf{Variabili Statiche/Campi vs Variabili Locali (\emph{Definite
  Assignment Analysis})}:

  \begin{itemize}
  \tightlist
  \item
    \passthrough{\lstinline!static float f1;!} e
    \passthrough{\lstinline!static String s1;!}: sono campi a livello di
    classe. \textbf{La JVM li inizializza SEMPRE automaticamente al loro
    default}: \passthrough{\lstinline!f1!} diventa
    \passthrough{\lstinline!0.0f!} e \passthrough{\lstinline!s1!}
    diventa \passthrough{\lstinline!null!}.
  \item
    \passthrough{\lstinline!int i2;!} e
    \passthrough{\lstinline!String s2;!}: sono variabili locali allocate
    nel frame dello \textbf{Stack}. \textbf{Le variabili locali NON
    hanno valori di default}.
  \item
    Se tentiamo di leggere \passthrough{\lstinline!i2!} o
    \passthrough{\lstinline!s2!} prima di aver assegnato loro un valore
    (\passthrough{\lstinline!System.out.println(i2)!}), il compilatore
    Java blocca la compilazione con l'errore:
    \passthrough{\lstinline!variable i2 might not have been initialized!}.
  \end{itemize}
\item
  \textbf{Suffisso dei Floating Point (\passthrough{\lstinline!3.4f!})}:

  \begin{itemize}
  \tightlist
  \item
    \passthrough{\lstinline!static float f2 = 3.4f;!}: il letterale
    \passthrough{\lstinline!3.4!} senza suffissi è un
    \passthrough{\lstinline!double!} a 64 bit. Scrivere
    \passthrough{\lstinline!float f = 3.4;!} genera errore di
    compilazione (\emph{loss of precision}). Il suffisso
    \passthrough{\lstinline!f!} forza il letterale a 32 bit.
  \end{itemize}
\item
  \textbf{Costanti con \passthrough{\lstinline!final!}}:

  \begin{itemize}
  \tightlist
  \item
    \passthrough{\lstinline!static final double PI = 3.14;!}: la keyword
    \passthrough{\lstinline!final!} rende la variabile a sola lettura
    dopo l'inizializzazione. Il tentativo di riassegnamento
    \passthrough{\lstinline!PI = 3.0;!} viene respinto dal compilatore
    (\passthrough{\lstinline!cannot assign a value to final variable PI!}).
    Per convenzione Java, le costanti
    \passthrough{\lstinline!static final!} si scrivono in
    \passthrough{\lstinline!MAIUSCOLO\_SNAKE\_CASE!}.
  \end{itemize}
\item
  \textbf{I 3 modi di rappresentare un \passthrough{\lstinline!char!}}:

  \begin{itemize}
  \tightlist
  \item
    \passthrough{\lstinline!c1 = 'V'!}: letterale carattere.
  \item
    \passthrough{\lstinline!c2 = 86!}: valore intero decimale
    ASCII/Unicode per la `V'.
  \item
    \passthrough{\lstinline!c3 = '\\u0056'!}: sequenza di escape Unicode
    a 16 bit in esadecimale (\(5 \times 16 + 6 = 86\)).
  \item
    Stampati a video, producono tutti e tre il carattere
    \passthrough{\lstinline!V!}.
  \end{itemize}
\item
  \textbf{Concatenazione di \passthrough{\lstinline!null!} vs
  \passthrough{\lstinline!NullPointerException!}}:

  \begin{itemize}
  \tightlist
  \item
    \passthrough{\lstinline!System.out.println("s1: " + s1);!} stampa
    \passthrough{\lstinline!"s1: null"!}. L'operatore
    \passthrough{\lstinline!+!} converte in sicurezza il riferimento
    nullo nella stringa letterale \passthrough{\lstinline!"null"!}.
  \item
    \passthrough{\lstinline!System.out.println(s1.toString());!} genera
    invece un \textbf{Run-Time Error}
    (\passthrough{\lstinline!NullPointerException!}): non è possibile
    dereferenziare un puntatore nullo per invocarvi un metodo d'istanza.
  \end{itemize}
\item
  \textbf{Operatori Logici Standard (\passthrough{\lstinline!|!},
  \passthrough{\lstinline!\&!}) vs Corto-Circuito (\emph{Short-Circuit}
  \passthrough{\lstinline!||!}, \passthrough{\lstinline!\&\&!}) e Side
  Effects}:

  \begin{itemize}
  \item
    Caso \passthrough{\lstinline!true | k++ == 0!}:

    \begin{itemize}
    \tightlist
    \item
      L'operatore singolo \passthrough{\lstinline!|!} (non a corto
      circuito) \textbf{valuta SEMPRE entrambi gli operandi}.
    \item
      Anche se la parte sinistra è già \passthrough{\lstinline!true!},
      la parte destra \passthrough{\lstinline!k++ == 0!} viene valutata:
      \passthrough{\lstinline!k!} viene confrontato con
      \passthrough{\lstinline!0!} (dando \passthrough{\lstinline!true!})
      e poi incrementato di \(1\).
    \item
      Stampa a video: \passthrough{\lstinline!true!}, seguito da
      \passthrough{\lstinline!k (standard eval): 1!}.
    \end{itemize}
  \item
    Caso \passthrough{\lstinline!true || k++ == 0!}:

    \begin{itemize}
    \tightlist
    \item
      L'operatore doppio \passthrough{\lstinline!||!} (a corto circuito)
      verifica la parte sinistra: poiché è
      \passthrough{\lstinline!true!}, il risultato globale dell'or
      logico è già matematicamente certo
      (\passthrough{\lstinline!true!}).
    \item
      La parte destra viene \textbf{completamente ignorata e mai
      valutata}!
    \item
      Di conseguenza, \passthrough{\lstinline!k++!} \textbf{non viene
      eseguito}.
    \item
      Stampa a video: \passthrough{\lstinline!true!}, seguito da
      \passthrough{\lstinline!k (short-circuit eval): 0!}.
    \end{itemize}
  \item
    \emph{Importanza pratica}: lo short-circuit
    \passthrough{\lstinline!\&\&!} è l'idioma fondamentale per evitare
    crash su \passthrough{\lstinline!null!}:

\begin{lstlisting}[language=Java]
if (s != null && s.length() > 0) { ... } // Se s è null, non chiama mai s.length()!
\end{lstlisting}
  \end{itemize}
\item
  \textbf{Scope delle Variabili e Divieto di Shadowing Locale}:

  \begin{itemize}
  \item
    In Java, il ciclo di vita e la visibilità delle variabili sono
    delimitati dalle graffe \passthrough{\lstinline!\{ ... \}!}.
  \item
    \passthrough{\lstinline!int n = 4; n++;!} dichiara
    \passthrough{\lstinline!n!} nello scope del metodo. Scrivere
    \passthrough{\lstinline!int n = 8;!} nello stesso scope o in un
    sottoblocco interno genera l'errore:
    \passthrough{\lstinline!Variable 'n' is already defined in the scope!}.
    In Java le variabili locali non possono oscurare altre variabili
    locali omonime.
  \item
    Blocchi anonimi indipendenti:

\begin{lstlisting}[language=Java]
{ int m = 7; m++; }
{ int m = 8; m--; }
\end{lstlisting}

    La variabile \passthrough{\lstinline!m!} del primo blocco viene
    distrutta alla chiusura della prima graffa. La seconda
    \passthrough{\lstinline!m!} è una nuova variabile in uno scope
    disgiunto. All'esterno dei blocchi, \passthrough{\lstinline!m!} non
    esiste (\passthrough{\lstinline!cannot find symbol variable m!}).
  \end{itemize}
\end{enumerate}

\begin{center}\rule{0.5\linewidth}{0.5pt}\end{center}

\subsubsection{3. Test: Output Esatto di
Esecuzione}\label{test-output-esatto-di-esecuzione}

Eseguendo \passthrough{\lstinline!Example.java!} da terminale o
IntelliJ:

\begin{lstlisting}
i1: 6
f1: 0.0
f2: 3.4
c1: V c2: V c3: V
s1: null
true
k (standard eval): 1
true
k (short-circuit eval): 0
n: 5
\end{lstlisting}

\begin{center}\rule{0.5\linewidth}{0.5pt}\end{center}

\begin{studynote}[Nota di Studio] Con la \textbf{Lezione 4 / T04} abbiamo chiarito le basi di
tipo, stack/heap e operatori.

Il prossimo blocco raggruppa le \textbf{Lezioni 05-06 / Slide T05
(\emph{Java Basic Syntax}) e T06 (\emph{Java Classes and Objects})},
dove nasce il progetto incrementale centrale del corso: - Struttura
della prima classe ad oggetti:
\href{file:///home/wally/Documenti/GitHub/Programmazione-II/25-26/Teoria-e-Esercitazioni/Codice-20260902/Lezioni05-06/SimpleDate}{\passthrough{\lstinline!SimpleDate!}}
(\passthrough{\lstinline!Date.java!} e
\passthrough{\lstinline!MainDate.java!}). - La classe sperimentale
\href{file:///home/wally/Documenti/GitHub/Programmazione-II/25-26/Teoria-e-Esercitazioni/Codice-20260902/Lezioni05-06/StringPlayground.java}{\passthrough{\lstinline!StringPlayground.java!}}.
- La prima validazione delle date: algoritmo dei giorni del mese e
calcolo degli anni bisestili secondo il calendario Gregoriano. -
Sintassi Java: costrutti di controllo
(\passthrough{\lstinline!if-else!}, \passthrough{\lstinline!switch!},
\passthrough{\lstinline!for!}, \passthrough{\lstinline!while!},
\passthrough{\lstinline!do-while!}), etichette (\emph{labeled
break/continue}). - Metodi, parametri, ritorno e la keyword
\passthrough{\lstinline!this!}. - La classe
\passthrough{\lstinline!MainDate!}: test delle istanze, instanziazione
con \passthrough{\lstinline!new!} e stampa dello stato.
\end{studynote}

\textbf{Dammi conferma per aprire le Lezioni 05-06 / T05-T06 e
analizzare la prima versione di \passthrough{\lstinline!SimpleDate!} e
\passthrough{\lstinline!StringPlayground!}!}


\end{document}
